% GENERATED FILE. Edit canonical lessons, then run npm run generate:books.
% canonical-source-hash: fnv1a64:d74852db68747e11
% canonical-lessons: ES-C561-marcharse, ES-C561-zarpar, ES-C561-partida, ES-C561-regreso, ES-C561-ausencia

\chapter{La partida}
\label{ch:es-la-partida}
\hlchaptermodality{561}

\begin{chapteropening}
\textbf{By the end of this chapter:} \emph{I can say that someone left, a ship set sail, someone was away and came back.}

Complete the last lesson of chapter 561: I can say that someone left, a ship set sail, someone was away and came back.
\end{chapteropening}

\section[marcharse]{marcharse — to leave, to go away}

\label{lesson:ES-C561-marcharse}



\begin{quote}
Before the new one: say the Spanish for the get-together, then the Spanish for the reunion.
\end{quote}

\subsection*{You'll want to know: marcharse}
\textbf{marcharse} — \textquotedblleft{}to leave, to go away\textquotedblright{}.

\begin{quote}
Se marchó sin decir adiós.
\end{quote}

\begin{quote}
\emph{He left without saying goodbye.}
\end{quote}

Common in Spain; \emph{irse} is used everywhere.

\begin{cousinweb}
\emph{marcharse} is built on \emph{marchar}, which was borrowed from French \emph{marcher} 'to walk, to march' (EtymDB). The French verb is usually traced to a Frankish word for marking out, pressing with the feet.
\end{cousinweb}

\begin{grammarlens}[title={What you've built}]
One more word for telling, explaining and describing at B1.
\end{grammarlens}

\subsection*{Your turn}
\begin{itemize}
\raggedright
  \item \emph{Say it:} \emph{marcharse}
  \item \emph{Say it:} \emph{marcharse}, once more
  \item \emph{Say it:} a sentence of your own with \emph{marcharse}
  \item \emph{Recall:} say the Spanish for the get-together, then the Spanish for the reunion, then say \emph{marcharse} again
\end{itemize}

\begin{tcolorbox}[breakable,colback=teal!4,colframe=teal!35!black,title={Before you move on}]
What does marcharse mean? (\textquotedblleft{}To leave, to go away\textquotedblright{}.) Say it once more.
\end{tcolorbox}
\section[zarpar]{zarpar — to set sail, to leave port}

\label{lesson:ES-C561-zarpar}



\begin{quote}
Before the new one: say the Spanish for the reunion, then the Spanish for to leave.
\end{quote}

\subsection*{You'll want to know: zarpar}
\textbf{zarpar} — \textquotedblleft{}to set sail, to leave port\textquotedblright{}.

\begin{quote}
El barco zarpó a las ocho.
\end{quote}

\begin{quote}
\emph{The ship set sail at eight.}
\end{quote}

Only for ships and boats.

\begin{cousinweb}
\emph{zarpar} comes from old Italian \emph{sarpare} 'to weigh anchor' (EtymDB). Where the Italian word came from is not settled.
\end{cousinweb}

\begin{grammarlens}[title={What you've built}]
One more word for telling, explaining and describing at B1.
\end{grammarlens}

\subsection*{Your turn}
\begin{itemize}
\raggedright
  \item \emph{Say it:} \emph{zarpar}
  \item \emph{Say it:} \emph{zarpar}, once more
  \item \emph{Say it:} a sentence of your own with \emph{zarpar}
  \item \emph{Recall:} say the Spanish for the reunion, then the Spanish for to leave, then say \emph{zarpar} again
\end{itemize}

\begin{tcolorbox}[breakable,colback=teal!4,colframe=teal!35!black,title={Before you move on}]
What does zarpar mean? (\textquotedblleft{}To set sail, to leave port\textquotedblright{}.) Say it once more.
\end{tcolorbox}
\section[la partida]{la partida — the departure; the game}

\label{lesson:ES-C561-partida}



\begin{quote}
Before the new one: say the Spanish for to leave, then the Spanish for to set sail.
\end{quote}

\subsection*{You'll want to know: la partida}
\textbf{la partida} — \textquotedblleft{}the departure; the game\textquotedblright{}.

\begin{quote}
La partida del tren se retrasó.
\end{quote}

\begin{quote}
\emph{The train's departure was delayed.}
\end{quote}

\emph{Una partida de cartas} is a game of cards. \emph{Un partido} is a match.

\begin{cousinweb}
\emph{partida} is built on the participle of \emph{partir}, which is inherited from Latin \emph{partīre} 'to divide' (EtymDB).
\end{cousinweb}

\begin{grammarlens}[title={What you've built}]
One more word for telling, explaining and describing at B1.
\end{grammarlens}

\subsection*{Your turn}
\begin{itemize}
\raggedright
  \item \emph{Say it:} \emph{la partida}
  \item \emph{Say it:} \emph{la partida}, once more
  \item \emph{Say it:} a sentence of your own with \emph{la partida}
  \item \emph{Recall:} say the Spanish for to leave, then the Spanish for to set sail, then say \emph{la partida} again
\end{itemize}

\begin{tcolorbox}[breakable,colback=teal!4,colframe=teal!35!black,title={Before you move on}]
What does la partida mean? (\textquotedblleft{}The departure; the game\textquotedblright{}.) Say it once more.
\end{tcolorbox}
\section[el regreso]{el regreso — the return}

\label{lesson:ES-C561-regreso}



\begin{quote}
Before the new one: say the Spanish for to set sail, then the Spanish for the departure.
\end{quote}

\subsection*{You'll want to know: el regreso}
\textbf{el regreso} — \textquotedblleft{}the return\textquotedblright{}.

\begin{quote}
A su regreso, todo había cambiado.
\end{quote}

\begin{quote}
\emph{On his return, everything had changed.}
\end{quote}

\emph{De regreso} — on the way back.

\begin{cousinweb}
\emph{regreso} is a learned borrowing from Latin \emph{regressus} 'a going back' (Lewis \& Short), from \emph{regredī} 'to step back'.
\end{cousinweb}

\begin{grammarlens}[title={What you've built}]
One more word for telling, explaining and describing at B1.
\end{grammarlens}

\subsection*{Your turn}
\begin{itemize}
\raggedright
  \item \emph{Say it:} \emph{el regreso}
  \item \emph{Say it:} \emph{el regreso}, once more
  \item \emph{Say it:} a sentence of your own with \emph{el regreso}
  \item \emph{Recall:} say the Spanish for to set sail, then the Spanish for the departure, then say \emph{el regreso} again
\end{itemize}

\begin{tcolorbox}[breakable,colback=teal!4,colframe=teal!35!black,title={Before you move on}]
What does el regreso mean? (\textquotedblleft{}The return\textquotedblright{}.) Say it once more.
\end{tcolorbox}
\section[la ausencia]{la ausencia — the absence}

\label{lesson:ES-C561-ausencia}



\begin{quote}
Before the new one: say the Spanish for the departure, then the Spanish for the return.
\end{quote}

\subsection*{You'll want to know: la ausencia}
\textbf{la ausencia} — \textquotedblleft{}the absence\textquotedblright{}.

\begin{quote}
Durante su ausencia, su hermana cuidó la casa.
\end{quote}

\begin{quote}
\emph{During his absence, his sister looked after the house.}
\end{quote}

\begin{cousinweb}
\emph{ausencia} is a learned borrowing from Latin \emph{absentia} 'absence' (Lewis \& Short), from \emph{abesse} 'to be away', built on \emph{ab} and \emph{esse} 'to be'.
\end{cousinweb}

\begin{grammarlens}[title={What you've built}]
One more word for telling, explaining and describing at B1.
\end{grammarlens}

\subsection*{Your turn}
\begin{itemize}
\raggedright
  \item \emph{Say it:} \emph{la ausencia}
  \item \emph{Say it:} \emph{la ausencia}, once more
  \item \emph{Say it:} a sentence of your own with \emph{la ausencia}
  \item \emph{Recall:} say the Spanish for the departure, then the Spanish for the return, then say \emph{la ausencia} again
\end{itemize}

\begin{tcolorbox}[breakable,colback=teal!4,colframe=teal!35!black,title={Before you move on}]
What does la ausencia mean? (\textquotedblleft{}The absence\textquotedblright{}.) Say it once more.
\end{tcolorbox}
